%%%%%%%%%%%%%%%%%%%%%%%%%%%%%%%%%%%%%%%%%%%%%%%%%%%%%%%%%%%%%%%%%%
%% Toto je inkludovaný LaTeXový soubor `ca_hyper.tex'.          %%
%% ------------------------------------------------------------ %%
%% Verze 2022-06-03                                             %%
%% Vyvíjí & udržuje <stanislav.hledik@physics.slu.cz>           %%
%%%%%%%%%%%%%%%%%%%%%%%%%%%%%%%%%%%%%%%%%%%%%%%%%%%%%%%%%%%%%%%%%%

\chapter{Konfigurace balíčku \pkg{hyperref}}\label{hyconf}%%%%%%%%%%%%%%%%%

Balíček \pkg{hyperref}~\cite{LtX:LIVE:CTAN:hyperref} je načten a
předkonfigurován již v~dokumentní třídě \pgm{\ipotname .cls} v~argumentu makra
\pgm{\textbackslash AtEndPreamble}~\cite{Wri-Leh:LIVE:CTAN:etoolbox}, aby byl
načten jako poslední ze všech balíčků, i~těch, jenž si uživatel vyžádá
v~preambuli hlavního souboru (viz též Dodatek~\ref{packages} na
stránce~\pageref{pkghyperref}).\footnote{To je z~technického hlediska žádoucí,
  vysvětlení lze nalézt v~manuálu
  \pkg{hyperref}~\cite{LtX:LIVE:CTAN:hyperref}.}  Načtení balíčku
\pkg{hyperref} se zapíná volbou \opt{htext} bez ohledu na to, která z~voleb
\opt{draft} či \opt{final} (případně i~ostatních popsaných
v~Dodatku~\ref{options}) je aktivní. Ovladač~\cite{SE:LIVE::DriversCreatePDF}
je autodetekován; obecně platí, že při kompilaci pomocí \pkg{pdflatex} se
použije odpovídající ovladač automaticky, při kompilaci pomocí \pkg{latex} se
použije ovladač \pkg{dvipdfmx} rovněž automaticky.\footnote{Toto chování lze
  změnit v~konfiguračním souboru \pgm{hyperref.cfg}. Ovladač \pkg{dvipdfmx}
  předpokládá následnou jím provedenou konverzi do
  \designation{pdf}. Praktický význam to však nemá, neboť kompilace do
  souborového formátu \designation{pdf} se téměř výhradně provádí pomocí
  \pkg{pdflatex}.}

Dodatečnou konfiguraci je možné provést pomocí specializovaného konfiguračního
souboru \pgm{hyperref.cfg} balíčku \pkg{hyperref} a volitelného konfiguračního
souboru \pgm{\ipotname .cfg} balíčku \pkg{\ipotname} (viz
Dodatek~\ref{config}) umístěných v~adresáři \pgm{manual}. Důvodem přesunu
určité malé části konfigurace do \pgm{\ipotname .cfg}~-- jmenovitě specifikace
údajů zobrazovaných v~menu \uv{Document Properties} Adobe Acrobat
Readeru~\cite{Adobe:LIVE::AcrobatReader} (viz obrázek~\ref{AcroReadDocProps})
i~jiných prohlížečů~\designation{pdf},
např.~M$\mu$PDF~\cite{Art:LIVE::MuPDF}~-- je zachování těchto údajů při
kompilaci pomocí \pkg{pdflatex} i~při vypnuté volbě \opt{htext}, typicky pro
tiskovou formu práce. Kdyby byly tyto údaje uvedeny pouze
v~\pgm{hyperref.cfg}, nenačetly by se. Nezapomeňte tyto údaje v~\pgm{\ipotname
  .cfg}, stejně jako hodnotu klíče \pgm{baseurl} v~\pgm{hyperref.cfg},
aktualizovat podle vašich potřeb. Další informace a inspiraci pro vhodné změny
lze najít v~Dodatku~\ref{config}.

Zkrácený obsah aktuálního konfiguračního souboru \pgm{hyperref.cfg} načtený
přímo z~něj (pomocí balíčku \pkg{fancyvrb}~\cite{Vos:LIVE:CTAN:fancyvrb})
vypadá následovně:

\VerbatimInput[frame=none,firstline=16,lastline=27,fontsize=\footnotesize]%
{hyperref.cfg}
\noindent$\langle$\emph{řádky vynechány}$\rangle$
\VerbatimInput[frame=none,firstline=39,lastline=41,fontsize=\footnotesize]%
{hyperref.cfg}
\noindent$\langle$\emph{řádky vynechány}$\rangle$
\VerbatimInput[frame=none,firstline=66,lastline=72,fontsize=\footnotesize]%
{hyperref.cfg}
\noindent$\langle$\emph{řádky vynechány}$\rangle$
\VerbatimInput[frame=none,firstline=85,lastline=89,fontsize=\footnotesize]%
{hyperref.cfg}

Balíček \pkg{hyperref} je poměrně složitý. Pro bližší informace o~něm a jeho
konfiguraci proto odkazujeme na rozsáhlou
dokumentaci~\cite{LtX:LIVE:CTAN:hyperref} nebo na některou
z~monografií~\cite{Goo-Mit-Sam:1994:Companion:,%
  Goo-Rah:1994:WebComp:,%
  Kop-Dal:2004:LaTeXPodrPruv:}.

Za zmínku stojí možnost nastavení klíčů \pgm{backref} a \pgm{pagebackref} na
hodnotu \pgm{true}, jež k~bibliografickým položkám přidá odkazy vedoucí zpět
na citace, z~nichž čtenář odskočil. Dále, klíč \pgm{pdfpagelayout} je
implicitně nastaven na hodnotu \pgm{TwoPageRight}, s~níž se v~Adobe Acrobat
Readeru \designation{pdf} dokument zobrazuje jako otevřená kniha s~lichou
stránkou vpravo, jak je ilustrováno na obrázku~\ref{AcroReadPgDisp}.

Definice příkazu \pgm{\textbackslash href} je rozšířena tak, aby byl
použitelný i~bez balíčku \pkg{hyperref}, tj. se zakomentovanou
volbou~\opt{htext}.

S~hypertextem souvisejí i~pole \pgm{eprint} a \pgm{doi} implementovaná
v~bibliografickém stylu \pgm{\ipotname .bst}; jejich hodnoty budou
hypertextovými odkazy na odpovídající bibliografickou položku
v~\designation{arXiv} a \designation{doi}. Podrobnější informace o~jejich
použití a tvorbě bibliografie obecně je v~Kapitole~\ref{refs}.

\begin{figure}[t]
\setlength{\fboxsep}{0sp}
\begin{center}
\fbox{\includegraphics[width=.99\linewidth]{AcroRead_DocProps}}
\end{center}
\caption[Vlastnosti dokumentu zobrazené Adobe Acrobat Readerem]%
{\label{AcroReadDocProps}Vlastnosti dokumentu zobrazené Adobe Acrobat Readerem.}
\end{figure}

\begin{figure}[t]
\setlength{\fboxsep}{0sp}
\begin{center}
\fbox{\includegraphics[width=.99\linewidth]{AcroRead_PgDisp_CZE}}
\end{center}
\caption[Zobrazení \uv{otevřené knihy} Adobe Acrobat Readerem]%
{\label{AcroReadPgDisp}Zobrazení \uv{otevřené knihy} Adobe Acrobat Readerem.}
\end{figure}

\endinput

%%
%% End of file `ca_hyper.tex'.
